\documentclass[10pt,a4paper]{amsart}
\usepackage[margin=19mm]{geometry}
\usepackage{amsmath,amssymb,mathtools}
\usepackage[scr=rsfs,cal=euler]{mathalfa}
\usepackage{xfrac}
\usepackage[unicode,hidelinks,bookmarks=false]{hyperref}
\newcommand{\ie}{i.e.\ }
\newcommand{\cf}{cf.\ }
\newcommand{\resp}{respectively\ }
\newcommand{\Subsection}{\subsection}
\providecommand{\nobreakdash}{\nobreak\hbox{-}}
\setlength{\emergencystretch}{2em}
\multlinegap0pt
\usepackage{amssymb}
\usepackage{bm}
\usepackage{mathtools}
\usepackage{dsfont}
\usepackage[normalem]{ulem}
\usepackage[shortlabels]{enumitem}
\usepackage{tikz}
\newtheorem{lemma}{Lemma}
\usepackage{cleveref}
\crefname{lemma}{Lemma}{Lemmas}
\newcommand{\iintss}{{\int\hspace{-0.28cm}\int}}
\newcommand{\la}{\lambda}
\newcommand{\miint}[1]{{\miints_{\hspace{-0.13cm}#1}}}
\newcommand{ \miints }{{\iintss\hspace{-0.56cm} -\hspace{-0.15cm}-}}
\title{Gradient higher integrability for degenerate parabolic double phase systems with two modulating coefficients}
\author{Jehan Oh, Abhrojyoti Sen}
\date{Source: arXiv:2504.15799v2 (CC BY 4.0)}
\begin{document}
\maketitle

\noindent\emph{Source and licence.} Gradient higher integrability for degenerate parabolic double phase systems with two modulating coefficients. Version arXiv:2504.15799v2 (CC BY 4.0), distributed by arXiv under the Creative Commons Attribution 4.0 International licence, http://creativecommons.org/licenses/by/4.0/, as recorded in the arXiv licence metadata for this exact version and shown on the abstract page https://arxiv.org/abs/2504.15799v2. Full licence notice: \texttt{licenses/arxiv-2504.15799-CC-BY-4.0.txt}.\par\smallskip

\noindent\textit{Editorial scope.} Counted unit: the article's lemma 3.13 (source0.tex lines 1563--1572). The article's complete original proof of it is delivered with it (source0.tex lines 1573--1698, one \texttt{proof} environment of 126 lines). No mathematics is written, repaired or re-derived: the statement and the proof are the article's own lines, in the article's own notation, and no retained line is substituted or re-flowed.  Retained uncounted as the source-local context the proof needs: lines 392--394; lines 947--952; lines 949--951; lines 1530--1539; lines 1541--1550.  Nothing was omitted from any retained span, so the package carries no omission record.  No retained line contains a citation, so the wrapper carries no bibliography and no bibliography entry.  No retained line contains a figure environment, an image-inclusion command or drawing code, so no image is delivered and the package declares no asset dependency.  Editorial formatting only, in addition: an AMS \texttt{amsart} preamble that restates the article's own declarations and macros used by the retained lines, namely the environments \texttt{lemma} on separate counters, together with the standard packages those lines need.  The retained environments print in this document as Lemma 1 in order of appearance, whereas the article prints the same items with the numbering of its own full document.  The complete native source of the article as distributed in the arXiv e-print for this exact version is bundled unchanged as \texttt{sources/176898/source0.tex}.
\begin{equation}    \label{eq: main equation}
    u_t -\operatorname*{div}\mathcal{A}(z,u,Du)=-\operatorname*{div}\mathcal{B}(z,F) \quad \text{in $\Omega_T$,}
\end{equation}

\begin{lemma}\label{lem : parabolic poincare inequality in Section 3}
    Let $u$ be a weak solution to \eqref{eq: main equation}. Then for $U_{R,S}(z_0)\subset \Omega_T$, $m\in(1,q]$ and $\theta\in(1/m,1]$, there exists a constant $c$ depending on $n$, $N$, $m$ and $L$ such that 
    \begin{align}\label{eq : parabolic poincare inequality in Section 3}
        \miint{U_{R,S}(z_0)} &\frac{|u-(u)_{U_{R,S}(z_0)}|^{\theta m}}{R^{\theta m}}\, dz \nonumber \\ &\leq c\miint{U_{R,S}(z_0)} |Du|^{\theta m} \,dz +c\left(\frac{S}{R^2}\miint{U_{R,S}(z_0)} [\widetilde{H}(z,|Du|)+\widetilde{H}(z,|F|)]\, dz \right)^{\theta m}.
    \end{align}
\end{lemma}

    \begin{align}
        \miint{U_{R,S}(z_0)} &\frac{|u-(u)_{U_{R,S}(z_0)}|^{\theta m}}{R^{\theta m}}\, dz \nonumber \\ &\leq c\miint{U_{R,S}(z_0)} |Du|^{\theta m} \,dz +c\left(\frac{S}{R^2}\miint{U_{R,S}(z_0)} [\widetilde{H}(z,|Du|)+\widetilde{H}(z,|F|)]\, dz \right)^{\theta m}.
    \end{align}

\begin{equation}\label{Eq 3.8}    
    \left\{
    \begin{aligned}
        &Ka(z_0)\lambda^p \leq b(z_0)\lambda^q,\\
        &\frac{a(z)}{3}\leq a(z_0)\leq 3a(z) \quad \text{and}\quad \frac{b(z)}{3}\leq b(z_0)\leq 3b(z) \ \, \text{for all}\,\,\, z\in G_{\rho, \la}(z_0),\\
        &\miint{G_{\tau,\lambda}(z_0)} [H(z,|Du|)+H(z,|F|)] \; dz <a(z_0)\lambda^p+b(z_0)\lambda^q \ \, \text{for every }\tau\in(\rho,c_v(2\rho)],\\
        &\miint{G_{\rho,\lambda}(z_0)} [H(z,|Du|)+H(z,|F|)] \; dz =a(z_0)\lambda^p+b(z_0)\lambda^q,
    \end{aligned}
    \right.
\end{equation}

\begin{equation}\label{Eq 3.9}    
    \left\{
    \begin{aligned}
        &Ka(z_0)\lambda^p \geq b(z_0)\lambda^q,\\
        &\frac{a(z)}{3}\leq a(z_0)\leq 3a(z) \quad \text{and}\quad \frac{b(z)}{3}\leq b(z_0)\leq 3b(z) \ \, \text{for all}\,\,\, z\in G_{\rho, \la}(z_0),\\
        &\miint{G_{\tau,\lambda}(z_0)} [H(z,|Du|)+H(z,|F|)] \; dz <a(z_0)\lambda^p+b(z_0)\lambda^q \ \, \text{for every }\tau\in(\rho,c_v(2\rho)],\\
        &\miint{G_{\rho,\lambda}(z_0)} [H(z,|Du|)+H(z,|F|)] \; dz =a(z_0)\lambda^p+b(z_0)\lambda^q.
    \end{aligned}
    \right.
\end{equation}

\begin{lemma}\label{Lemma 3.13}
Let $u$ be a weak solution to \eqref{eq: main equation}. Then there exists a constant $c=c(\textnormal{\texttt{data}}) \geq 1$ such that for any $G_{4\rho,\lambda}(z_0)\subset \Omega_T$ with \eqref{Eq 3.8}--\eqref{Eq 3.9}, $\tau \in [2\rho,4\rho]$ and $\theta \in ((q-1)/p,1]$, we have
\begin{align}\label{Eq 3.10}
&\miint{G_{\tau,\lambda}(z_0)} \frac{|u-(u)^{(p,q)}_{z_0;\tau,\lambda}|^{\theta p}}{\tau^{\theta p}}\; dz \leq c \miint{G_{\tau,\lambda}(z_0)} \left(|Du|+|F|\right)^{\theta p}\, dz
\end{align}
and 
\begin{align}\label{Eq: 3.11}
&\miint{G_{\tau,\lambda}(z_0)} \frac{|u-(u)^{(p,q)}_{z_0;\tau,\lambda}|^{\theta q}}{\tau^{\theta q}}\; dz \leq c \miint{G_{\tau,\lambda}(z_0)} \left(|Du|+|F|\right)^{\theta q}\, dz.
\end{align}
\end{lemma}

\begin{proof}
We observe from \eqref{eq : parabolic poincare inequality in Section 3} of \cref{lem : parabolic poincare inequality in Section 3} that
\begin{align}\label{EQQ3.12}
    &\miint{G_{\tau, \lambda}(z_0)}\frac{|u-(u)^{(p,q)}_{z_0, \tau, \lambda}|^{\theta p}}{\tau^{\theta p}}\,dz\nonumber\\
    &\leq c\miint{G_{\tau, \lambda}(z_0)}|Du|^{\theta p}\, dz \nonumber\\
    &\quad + c\left(\frac{\lambda^2}{a(z_0)\lambda^p+b(z_0)\lambda^q}\miint{G_{\tau, \lambda}(z_0)}a(z)\left(|Du|^{p-1}+|F|^{p-1}\right)+b(z)\left(|Du|^{q-1}+|F|^{q-1}\right) dz\right)^{\theta p}\nonumber\\
    &\leq c\miint{G_{\tau, \lambda}(z_0)}|Du|^{\theta p}\, dz+ c\underbrace{\left(\frac{\lambda^2}{a(z_0)\lambda^p+b(z_0)\lambda^q}\miint{G_{\tau, \lambda}(z_0)}a(z)\left(|Du|+|F|\right)^{p-1} dz\right)^{\theta p}}_{\mathrm{I}}\nonumber\\
    &\quad +c\underbrace{\left(\frac{\lambda^2}{a(z_0)\lambda^p+b(z_0)\lambda^q}\miint{G_{\tau, \lambda}(z_0)}b(z)\left(|Du|+|F|\right)^{q-1} dz\right)^{\theta p}}_{\mathrm{II}}.
\end{align}
\textbf{Case I:} In this case, we estimate $\mathrm{I}$ and $\mathrm{II}$ of \eqref{EQQ3.12} considering the assumption \eqref{Eq 3.8}. First we observe, since $b(z_0)\neq 0,$ from the first, second and the third conditions of \eqref{Eq 3.8}, we get
\begin{align*}
 \miint{G_{\tau, \lambda}(z_0)} b(z_0)\left(|Du|+|F|\right)^q\,dz \leq c \miint{G_{\tau, \lambda}(z_0)} b(z)\left(|Du|+|F|\right)^q\,dz&\leq c \left(a(z_0)\lambda^p+b(z_0)\lambda^q\right)\\
 &\leq c\, b(z_0)\lambda^q.
\end{align*}
Thus, we have
\begin{align}\label{bound by q}
    \miint{G_{\tau, \lambda}(z_0)}\left(|Du|+|F|\right)^q\, dz \leq c\lambda^q.
\end{align}

\noindent\textbf{Estimate of $\mathrm{I}$:} Using the second condition of \eqref{Eq 3.8} and H\"{o}lder's inequality, we obtain
\begin{align*}
    &\left(\frac{\lambda^2}{a(z_0)\lambda^p+b(z_0)\lambda^q}\miint{G_{\tau, \lambda}(z_0)}a(z)\left(|Du|+|F|\right)^{p-1}\,dz\right)^{\theta p}\\
    &\leq \left(\frac{\lambda^2}{a(z_0)\lambda^p+b(z_0)\lambda^q}\right)^{\theta p}a(z_0)^{\theta}\left(\miint{G_{\tau, \lambda}(z_0)}a(z)^{\frac{p-1}{p}}\left(|Du|+|F|\right)^{p-1}\,dz\right)^{\theta p}\\
    &\leq \left(\frac{\lambda^2}{a(z_0)\lambda^p+b(z_0)\lambda^q}\right)^{\theta p}a(z_0)^{\theta}\left(\miint{G_{\tau, \lambda}(z_0)}a(z)^\theta \left(|Du|+|F|\right)^{\theta p}\,dz\right)^{p-1}\\
    &\leq c\left(\frac{\lambda^2}{a(z_0)\lambda^p+b(z_0)\lambda^q}\right)^{\theta p}(a(z_0)\lambda^p+b(z_0)\lambda^q)^{\theta (p-2)}a(z_0)^{2\theta} \miint{G_{\tau, \lambda}(z_0)} \left(|Du|+|F|\right)^{\theta p}\,dz.
\end{align*}
Now using $a(z_0)\leq \frac{1}{K} b(z_0)\lambda^{q-p}$ from the assumption \eqref{Eq 3.8}, we have
\begin{align*}
 &\left(\frac{\lambda^2}{a(z_0)\lambda^p+b(z_0)\lambda^q}\right)^{\theta p}(a(z_0)\lambda^p+b(z_0)\lambda^q)^{\theta (p-2)}a(z_0)^{2\theta}\\
 &\quad \leq \left(\frac{\lambda^p}{a(z_0)\lambda^p+b(z_0)\lambda^q}\right)^{2\theta}b(z_0)^{2\theta}\lambda^{2\theta(q-p)}=\left(\frac{b(z_0)\lambda^{q}}{a(z_0)\lambda^p+b(z_0)\lambda^q}\right)^{2\theta} \leq 1.
\end{align*}
Thus we finally obtain
\begin{align*}
\left(\frac{\lambda^2}{a(z_0)\lambda^p+b(z_0)\lambda^q}\miint{G_{\tau, \lambda}(z_0)}a(z)\left(|Du|+|F|\right)^{p-1}\,dz\right)^{\theta p} \leq c \miint{G_{\tau, \lambda}(z_0)} \left(|Du|+|F|\right)^{\theta p}\,dz.    
\end{align*}

\noindent\textbf{Estimate of $\mathrm{II}$:} Similarly, using the second condition of \eqref{Eq 3.8} and H\"{o}lder's inequality, we get
\begin{align*}
&\left(\frac{\lambda^2}{a(z_0)\lambda^p+b(z_0)\lambda^q}\miint{G_{\tau, \lambda}(z_0)}b(z)\left(|Du|+|F|\right)^{q-1}\,dz\right)^{\theta p}\\
&\leq \left(\frac{\lambda^2 b(z_0)}{a(z_0)\lambda^p+b(z_0)\lambda^q}\right)^{\theta p}\left(\miint{G_{\tau, \lambda}(z_0)}\left(|Du|+|F|\right)^q\, dz\right)^{\frac{(q-2)\theta p}{q}} \miint{G_{\tau, \lambda}(z_0)}\left(|Du|+|F|\right)^{\theta p}\, dz\\
&\leq \left(\frac{\lambda^2 b(z_0)}{a(z_0)\lambda^p+b(z_0)\lambda^q}\right)^{\theta p} \lambda^{(q-2)\theta p} \miint{G_{\tau, \lambda}(z_0)}\left(|Du|+|F|\right)^{\theta p}\, dz\\
&\leq \miint{G_{\tau, \lambda}(z_0)}\left(|Du|+|F|\right)^{\theta p}\, dz.
\end{align*}
Thus, combining the above estimates, we obtain \eqref{Eq 3.10}.

\noindent\textbf{Case II:} In this case, we estimate $\mathrm{I}$ and $\mathrm{II}$ of \eqref{EQQ3.12} considering the assumption \eqref{Eq 3.9}. Since $a(z_0)\neq 0,$ using the first, second and the third assumptions of \eqref{Eq 3.9}, we get
\begin{align*}
    a(z_0)\miint{G_{\tau, \lambda}(z_0)}\left(|Du|+|F|\right)^p\,dz\leq 3 \miint{G_{\tau, \lambda}(z_0)}a(z)\left(|Du|+|F|\right)^p\,dz&\leq a(z_0)\lambda^p+b(z_0)\lambda^q\\&\leq c a(z_0)\lambda^p,
\end{align*}
and thus we conclude that
\begin{align}\label{bound by lambda-p}
    \miint{G_{\tau, \lambda}(z_0)}\left(|Du|+|F|\right)^p\,dz \leq c \lambda^p.
\end{align}
The estimates in this case are similar to those in Case I.

\noindent \textbf{Estimate of $\mathrm{I}$:} Using the second condition of \eqref{Eq 3.9}, \eqref{bound by lambda-p} and H\"{o}lder's inequality, we obtain
\begin{align*}
    &\left(\frac{\lambda^2}{a(z_0)\lambda^p+b(z_0)\lambda^q}\miint{G_{\tau, \lambda}(z_0)}a(z)\left(|Du|+|F|\right)^{p-1}\,dz\right)^{\theta p}\\
    & \leq c \left(\frac{a(z_0)\lambda^2}{a(z_0)\lambda^p+b(z_0)\lambda^q}\right)^{\theta p}\lambda^{\theta p (p-2)} \miint{G_{\tau, \lambda}(z_0)}\left(|Du|+|F|\right)^{\theta p}\, dz\\
    &=c\left(\frac{a(z_0)\lambda^p}{a(z_0)\lambda^p+b(z_0)\lambda^q}\right)^{\theta p} \miint{G_{\tau, \lambda}(z_0)}\left(|Du|+|F|\right)^{\theta p}\, dz \leq c\miint{G_{\tau, \lambda}(z_0)}\left(|Du|+|F|\right)^{\theta p}\, dz.
\end{align*}
\textbf{Estimate of $\mathrm{II}$:} Similarly, we estimate
\begin{align*}
 &\left(\frac{\lambda^2}{a(z_0)\lambda^p+b(z_0)\lambda^q}\miint{G_{\tau, \lambda}(z_0)}b(z)\left(|Du|+|F|\right)^{q-1}\,dz\right)^{\theta p}\\
 &\leq c \left(\frac{\lambda^2 b(z_0)}{a(z_0)\lambda^p+b(z_0)\lambda^q}\right)^{\theta p}\left(\miint{G_{\tau, \lambda}(z_0)}\left(|Du|+|F|\right)^{p}\,dz\right)^{\theta (q-2)} \miint{G_{\tau, \lambda}(z_0)}\left(|Du|+|F|\right)^{\theta p}\,dz\\
 &\leq c \left(\frac{\lambda^q b(z_0)}{a(z_0)\lambda^p+b(z_0)\lambda^q}\right)^{\theta p} \miint{G_{\tau, \lambda}(z_0)}\left(|Du|+|F|\right)^{\theta p}\,dz \leq \miint{G_{\tau, \lambda}(z_0)}\left(|Du|+|F|\right)^{\theta p}\,dz.
\end{align*}
Thus, combining the above estimates, we obtain \eqref{Eq 3.10}.

Next we prove the estimate \eqref{Eq: 3.11}. Again, it follows from \eqref{eq : parabolic poincare inequality in Section 3} of \cref{lem : parabolic poincare inequality in Section 3} that
\begin{align}\label{EQQ3.14}
&\miint{G_{\tau, \lambda}(z_0)}\frac{|u-(u)^{(p,q)}_{z_0, \tau, \lambda}|^{\theta q}}{\tau^{\theta q}}\,dz\nonumber\\
&\quad \leq c\miint{G_{\tau, \lambda}(z_0)}|Du|^{\theta q}\, dz+ c\underbrace{\left(\frac{\lambda^2}{a(z_0)\lambda^p+b(z_0)\lambda^q}\miint{G_{\tau, \lambda}(z_0)}a(z)\left(|Du|+|F|\right)^{p-1}\,dz\right)^{\theta q}}_{\mathrm{I}}\nonumber\\
&\qquad +c\underbrace{\left(\frac{\lambda^2}{a(z_0)\lambda^p+b(z_0)\lambda^q}\miint{G_{\tau, \lambda}(z_0)}b(z)\left(|Du|+|F|\right)^{q-1}\,dz\right)^{\theta q}}_{\mathrm{II}}.
\end{align}

\noindent \textbf{Case I:} In this case, we estimate $\mathrm{I}$ and $\mathrm{II}$ of \eqref{EQQ3.14} considering the assumption \eqref{Eq 3.8}.

\noindent \textbf{Estimate of $\mathrm{I}$:} Using the second condition of \eqref{Eq 3.8}, H\"{o}lder's inequality and \eqref{bound by q}, we see that
\begin{align*}
&\left(\frac{\lambda^2}{a(z_0)\lambda^p+b(z_0)\lambda^q}\miint{G_{\tau, \lambda}(z_0)}a(z)\left(|Du|+|F|\right)^{p-1}\,dz\right)^{\theta q}\\
&\leq c\left(\frac{\lambda^2 a(z_0)}{a(z_0)\lambda^p+b(z_0)\lambda^q}\right)^{\theta q}\left(\miint{G_{\tau, \lambda}(z_0)}\left(|Du|+|F|\right)^{\theta q}\,dz\right)^{p-1}\\
&\leq c \left(\frac{\lambda^2 a(z_0)}{a(z_0)\lambda^p+b(z_0)\lambda^q}\right)^{\theta q}\left(\miint{G_{\tau, \lambda}(z_0)}\left(|Du|+|F|\right)^q\, dz\right)^{\theta(p-2)} \miint{G_{\tau, \lambda}(z_0)}\left(|Du|+|F|\right)^{\theta q}\, dz\\
&\leq c\left(\frac{a(z_0)\lambda^p}{a(z_0)\lambda^p+b(z_0)\lambda^q}\right)^{\theta q} \miint{G_{\tau, \lambda}(z_0)}\left(|Du|+|F|\right)^{\theta q}\, dz \\
&\leq c \miint{G_{\tau, \lambda}(z_0)}\left(|Du|+|F|\right)^{\theta q}\, dz.
\end{align*}
\textbf{Estimate of $\mathrm{II}$:} Similarly, we can estimate
\begin{align*}
&\left(\frac{\lambda^2}{a(z_0)\lambda^p+b(z_0)\lambda^q}\miint{G_{\tau, \lambda}(z_0)}b(z)\left(|Du|+|F|\right)^{q-1}\,dz\right)^{\theta q}\\
&\leq c\left(\frac{\lambda^2 b(z_0)}{a(z_0)\lambda^p+b(z_0)\lambda^q}\right)^{\theta q}\left(\miint{G_{\tau, \lambda}(z_0)}\left(|Du|+|F|\right)^q\,dz\right)^{\theta (q-2)} \miint{G_{\tau, \lambda}(z_0)}\left(|Du|+|F|\right)^{\theta q}\,dz\\
&\leq c \left(\frac{\lambda^q b(z_0)}{a(z_0)\lambda^p+b(z_0)\lambda^q}\right)^{\theta q} \miint{G_{\tau, \lambda}(z_0)}\left(|Du|+|F|\right)^{\theta q}\,dz\\
&\leq c \miint{G_{\tau, \lambda}(z_0)}\left(|Du|+|F|\right)^{\theta q}\,dz.
\end{align*}
Combining the above estimates completes the proof of \eqref{Eq: 3.11}.



\noindent \textbf{Case II:} We estimate $\mathrm{I}$ and $\mathrm{II}$ of \eqref{EQQ3.14} by considering the assumption \eqref{Eq 3.9}.

\noindent\textbf{Estimate of $\mathrm{I}$:} Using \eqref{bound by lambda-p}, H\"{o}lder's inequality and the second condition of \eqref{Eq 3.9}, we get
\begin{align*}
&\left(\frac{\lambda^2}{a(z_0)\lambda^p+b(z_0)\lambda^q}\miint{G_{\tau, \lambda}(z_0)}a(z)\left(|Du|+|F|\right)^{p-1}\,dz\right)^{\theta q}\\
&\leq c\left(\frac{a(z_0)\lambda^2}{a(z_0)\lambda^p+b(z_0)\lambda^q}\right)^{\theta q}\left(\miint{G_{\tau, \lambda}(z_0)}\left(|Du|+|F|\right)^{\theta p}\,dz\right)^{\frac{(p-1)q}{p}}\\
&=c\left(\frac{a(z_0)\lambda^2}{a(z_0)\lambda^p+b(z_0)\lambda^q}\,\right)^{\theta q}\left(\miint{G_{\tau, \lambda}(z_0)}\left(|Du|+|F|\right)^{\theta p}\,dz\right)^{\frac{(p-2)q}{p}}\left(\miint{G_{\tau, \lambda}(z_0)}\left(|Du|+|F|\right)^{\theta p}\,dz\right)^{\frac{q}{p}}\\
&\leq c\left(\frac{a(z_0)\lambda^2}{a(z_0)\lambda^p+b(z_0)\lambda^q}\right)^{\theta q}\lambda^{\theta q(p-2)} \miint{G_{\tau, \lambda}(z_0)}\left(|Du|+|F|\right)^{\theta q} dz\\
&=c\left(\frac{a(z_0)\lambda^p}{a(z_0)\lambda^p+b(z_0)\lambda^q}\right) \miint{G_{\tau, \lambda}(z_0)}\left(|Du|+|F|\right)^{\theta q}\,dz\\
&\leq \miint{G_{\tau, \lambda}(z_0)}\left(|Du|+|F|\right)^{\theta q} dz.
\end{align*}

\noindent \textbf{Estimate of $\mathrm{II}$:} Similarly, we estimate
\begin{align*}
&\left(\frac{\lambda^2}{a(z_0)\lambda^p+b(z_0)\lambda^q}\miint{G_{\tau, \lambda}(z_0)}b(z)\left(|Du|+|F|\right)^{q-1}\,dz\right)^{\theta q}\\
&\leq c\left(\frac{\lambda^2}{a(z_0)\lambda^p+b(z_0)\lambda^q}\right)^{\theta q}b(z_0)^{\theta}\left(\miint{G_{\tau, \lambda}(z_0)}b(z)^{\frac{q-1}{q}}\left(|Du|+|F|\right)^{q-1}\, dz\right)^{\theta q}\\
&\leq c \left(\frac{\lambda^2}{a(z_0)\lambda^p+b(z_0)\lambda^q}\right)^{\theta q}b(z_0)^{\theta}\left(\miint{G_{\tau, \lambda}(z_0)}b(z)^{\theta}\left(|Du|+|F|\right)^{\theta q}dz\right)^{q-1}\\
&\leq c \left(\frac{\lambda^2}{a(z_0)\lambda^p+b(z_0)\lambda^q}\right)^{\theta q}(a(z_0)\lambda^p+b(z_0)\lambda^q)^{\theta (q-2)}b(z_0)^{2\theta} \miint{G_{\tau, \lambda}(z_0)} \left(|Du|+|F|\right)^{\theta q}dz.
\end{align*}
As in the previous case, using the first condition of \eqref{Eq 3.9}, i.e., $Ka(z_0)\lambda^p\geq b(z_0)\lambda^q,$ we obtain
\begin{align*}
 \left(\frac{\lambda^2}{a(z_0)\lambda^p+b(z_0)\lambda^q}\right)^{\theta q}(a(z_0)\lambda^p+b(z_0)\lambda^q)^{\theta (q-2)}b(z_0)^{2\theta} \leq 1.   
\end{align*}
It follows that
\begin{align*}
    \left(\frac{\lambda^2}{a(z_0)\lambda^p+b(z_0)\lambda^q}\miint{G_{\tau, \lambda}(z_0)}b(z)\left(|Du|+|F|\right)^{q-1}\,dz\right)^{\theta q} \leq c \miint{G_{\tau, \lambda}(z_0)} \left(|Du|+|F|\right)^{\theta q} dz.
\end{align*}
Thus, combining the above estimates, we obtain \eqref{Eq: 3.11} for this case.
\end{proof}

\end{document}
